\scanpage{84}
\uline{Sites fibrés}

\diagram{\begin{tikzcd}[ampersand replacement=\&,row sep=tiny]
\unsure{E|\mathcal E}\arrow[d]\\ B
\end{tikzcd}}

catégorie fibrée, avec les \deleted{foncteurs} \inserted{catégories} fibres
des sites, les foncteurs $f^*$ des morphismes
de sites [continus \unsure{suffira|suffit|suffisantes} pour certaines
choses). On met sur $E$ la top. la moins fine de
celles qui rend \unsure{continu|continus} les foncteurs $E_b\longrightarrow E$.

Les familles couvrantes sont décrites par VI 3.2.

Supposons $B,E\in\mathcal U$, \deleted{\unsure{tous et}} (pour simplifier) les
top. des $E_b$ \unsure{standart|A-endroit|Abondant}. Plongeons chaque
$E_b$ dans $\widetilde E_b$, on trouve une nouvelle catégorie
fibrée $\overline E$, : fibres les \inserted{topos} $\widetilde E_b$ ; un $f$-\deleted{\unsure{foncteur}} morphisme de
$F\in\unsure{\overline E_x|\widetilde E_x}$ dans $G\in\widetilde E_y$ est un morphisme

\margin{$\unsure{f}\quad x\longrightarrow y$}

\[
\text{dans }\widetilde E_x \qquad F\longrightarrow f^*(G)
\]

[en désignant aussi par $f^*$ l'ext.
de $f^*$ aux $\widetilde E_b$].

Déterminons le topos $\widetilde E$, \unsure{le|des|dans} \deleted{\unsure{uns|un}}
\unsure{dernier point|énoncé général} nous dit que \unsure{ses} \deleted{\unsure{sections}} \inserted{éléments} s'identifient

\[
\deleted{\widetilde E\simeq \underset{X\in\unsure{\mathrm{Ob}}\,E}{\underleftarrow{\lim}}
(E_{p(X)})/X = \overset{\text{\unsure{des fibres|des lieux}}}{\underset{X\in\mathrm{Ob}\, B}{\underleftarrow{\lim}}}E_x}
\]

\deleted{mais étant entendu} \unsure{,}
aux sections sur $E$ de la catégorie fibrée dont
la fibre, en $X\in\mathrm{Ob}\,E$, est la \deleted{\unsure{fibre}} \inserted{catégorie \uline{opposée} au} topos
associé au ``petit site'' défini par $X$ dans $E$, savoir

% archive sheet number: 83
